\section{Experimental platform}\label{sec:platform}
\subsection{Aircraft and autopilot}
The aircraft (Fig.~\ref{fig:aircraft}) is an \SI{8}{kg} twin-boom quadplane with a \SI{2}{m} foam-core wing of Selig S1223 section~\cite{Selig1997} (chord \SI{0.292}{m}, aspect ratio \numAR{}), four brushless lift motors at the boom corners, an H-tail and a tractor cruise motor with a folding propeller, flown by a Pixhawk-class flight controller running PX4~\cite{Meier2015}. The autopilot's estimator provides airspeed, flight-path angle, pitch and their rates at \SI{100}{Hz}; its multicopter attitude and rate loops stay active through the hand-off and receive the pitch setpoint and the rotor collective as an attitude target, while the pusher receives its command through an actuator-set message.
\begin{figure*}[t]\centering\includegraphics[width=\textwidth]{fig04_hardware.pdf}
\caption{The quadplane: airframe, power and autopilot bay, ground checks, hover on the lift rotors, cruise on the wing with the rotors stopped, cruise on the pusher, top view over the field, and landing.}\label{fig:aircraft}
\end{figure*}

\subsection{Identified model and parameter box}
The parameters of \eqref{eq:model} were identified from the PX4 log of a flight with several hand-offs, released as a dataset~\cite{Dony2026data}: the polar $(C_{L0},C_{L\alpha},C_{D0},k)$ from the lift and drag balances of the cruise segments anchored to the computational-fluid-dynamics polar of the wing, the attitude-loop rate and gain $(\kappa,\gth)$ from the pitch response to the setpoint, the thrust maps $(k_{\mathrm c},k_{\mathrm r})$ from the throttle-to-acceleration balances and the actuator lags from the transitions. Table~\ref{tab:box} lists the nominal values and the intervals that define the box $\Theta$: two standard deviations of the identification for the polar, the range of attitude-loop settings of the autopilot for $\kappa$ and $\gth$, and $\pm50\%$ of the lags and the identified spread of the thrust maps for the actuators. The power of the rotors follows the momentum-theory law $P_{\mathrm r}=c_{\mathrm r}T_{\mathrm r}^{3/2}$ through the measured hover power and the pusher power the linear law $P_{\mathrm c}=c_{\mathrm c}T_{\mathrm c}$ through the measured cruise power.
\begin{table}[t]\centering\footnotesize
\caption{Nominal parameters and the box $\Theta$ identified from the flight log.}\label{tab:box}
\setlength{\tabcolsep}{2.5pt}\scriptsize
\begin{tabular}{lcc|lcc}\toprule
parameter & nominal & interval & parameter & nominal & interval\\\midrule
$C_{L0}$ & 1.141 & $\pm0.144$ & $\kappa$ [1/s] & 2.0 & $[0.67,5.0]$\\
$C_{L\alpha}$ [1/rad] & 3.12 & $\pm25\%$ & $\gth$ & 1.0 & $[0.9,1.0]$\\
$C_{D0}$ & 0.133 & $\pm0.044$ & $\lambda_{\mathrm c}$ [1/s] & 2.0 & $[1.33,4.0]$\\
$k$ & 0.060 & $\pm25\%$ & $\lambda_{\mathrm r}$ [1/s] & 6.67 & $[4.44,13.3]$\\
$\overline T_{\mathrm c}$ [N] & 30.3 & $\beta_{\mathrm c}\in[0.78,1.22]$ & $\overline T_{\mathrm r}$ [N] & 118.7 & $\beta_{\mathrm r}\in[0.85,1.15]$\\
\bottomrule\end{tabular}\end{table}
The constraints are $\bar\alpha=\SI{5}{\degree}$ (inside the range verified unstalled in cruise), $\theta\in[-15,20]^\circ$, $\gamma\in[-15,15]^\circ$, $V_{\min}=\SI{4}{m/s}$, $K_L=\numKL{}$, $\dot\gamma_{\lim}=\SI{\numGdLim}{rad/s}$, the input box $|\thsp|\le\SI{30}{\degree}$, $T^{\mathrm c}_{\mathrm c}\in[0,\SI{30.3}{N}]$, $T^{\mathrm c}_{\mathrm r}\in[0,\SI{118.7}{N}]$ and the step bounds $\bar\delta=(\SI{4}{\degree},\SI{0.75}{N},\SI{4}{N})$ per \SI{50}{ms}. The corridor runs from the hover-exit airspeed $\rho_0=\SI{6}{m/s}$ to the cruise airspeed $\rho_{\mathrm c}=\SI{15}{m/s}$.

\subsection{Simulator, training and the safety module}\label{sec:num}
The differentiable model integrates \eqref{eq:model} with a Runge--Kutta step of two substeps per sample in PyTorch~\cite{Paszke2019} and the linear polar; the evaluation plant uses five substeps in NumPy and folds the lift curve beyond the stall angles of $10^\circ$ and $-14^\circ$ with a post-stall slope of \SI{-2}{/rad} and a drag rise, so that a policy that leaves the linear range is punished by the plant rather than rewarded by the model. HOT has $L=\numL{}$ tokens of width $d=\numD{}$, $N=\numBlocks{}$ blocks of \numHeads{} heads, \numParams{} parameters, and is trained for \numIters{} iterations of batch \numBatch{} with Adam, learning rate \numLr{} with cosine decay, gradient norm clipped at $1$, weights $(w_E,w_B,w_A,w_S)=(\numWE,\numWB,\numWA,\numWS)$ in \eqref{eq:loss}, gust standard deviation \SI{1}{m/s} and the noise of Table~\ref{tab:noise} on the window; the curriculum grows the horizon from \SI{3}{s} to \SI{15}{s} over the first $60\%$ of the iterations. The lift-safety module uses the plant set of the \numPlantsOnline{} vertices of the fast parameters and the lift-curve offset with \numSamplesOnline{} interior samples online, the backup gains $k_\gamma=\numKgam$, $k_{\mathrm r}=\numKr$~N per rad/s per step, the reference band $[\numRlo,\numRhi]$~rad/s, the lead $\tau_{\mathrm b}=\numTauB$~s, the pusher loop $k_{\mathrm c}=\numKc$~N per m/s$^2$ per step, $k_V=\numKv$~/s, $\bar a=\numAmax$~m/s$^2$ on the corridor $[\numVblo,\numVbhi]$~m/s, the horizon $H=\numH{}$ samples, the margins $m=(\numMargins)$, the settled neighbourhood $\mathcal C$: lift fraction in $[\numClfLo,\numClfHi]$, $|\gamma|\le\numCgam^\circ$, $|\dot\gamma|\le\numCgd$~rad/s, $\alpha\in[\numCalLo,\numCalHi]^\circ$, $V\in[\numCvLo,\numCvHi]$~m/s and $|\dot V|\le\numCvd$~m/s$^2$, and the measurement bound $\varepsilon=(\numEps)$. Assumption~\ref{ass:settle} was verified over the full plant set of \numPlantsFull{} plants; the one-step gain of Lemma~\ref{lem:corner} is \numKrGain{}.
\begin{table}[t]\centering\footnotesize
\caption{Sensor-noise standard deviations applied to the window during training and in the noise cell.}\label{tab:noise}
\begin{tabular}{lccccc}\toprule
& $V$ [m/s] & $\gamma$ [deg] & $\theta$ [deg] & $T_{\mathrm c}$ [N] & $T_{\mathrm r}$ [N]\\\midrule
level 1 & 0.3 & 0.5 & 0.3 & 0.6 & 2.4\\
\bottomrule\end{tabular}\end{table}

\subsection{PX4 firmware in the loop}
The autopilot firmware (PX4 v1.17) runs in software-in-the-loop with its simulation-in-hardware plant configured to the identified aircraft: mass, inertia, rotor thrust and lag, the wing with the polar of Table~\ref{tab:box}, the tail and the pusher lag. The firmware's estimator, attitude and rate loops, control allocation and offboard interface are unchanged; HOT runs as an offboard node that receives the states over MAVLink at \SI{100}{Hz}, sends the pitch setpoint and the rotor collective as an attitude target and the pusher command as an actuator-set command, and re-plans every \SI{50}{ms}. The network evaluates in \hotEvalMs{}~ms and the safety check takes \hotCheckMs{}~ms on average (\hotCheckMaxMs{}~ms at most) in interpreted code on one core, so the loop runs at real time on average and the simulation is slowed to half speed to keep every update inside its period.
